

\section{KS--Regularization} \label{ch:regularization}

The fourth main feature of the codes since NBODY3 is a special
treatment of close binaries.
A close encounter is characterised by an impact parameter
that is smaller than the parameter for a 90 degree deflection
%
\begin{equation}
p_{90} = 2 G(m_1 + m_2)/v_\infty^2   \label{eq:p90}
\end{equation}
%
where $G$, $m_1$, $m_2$, $v_\infty$ are the gravitational constant,
the masses of the two particles and their relative velocity
at infinity.
%In $N$--body units (Ch.~\ref{ch:nbtime}), it is
%$p_{90} \propto 1/N $ for a system in virial equilibrium
%(this is why these units are more important in small--$N$ systems).
%However, no matter how large $N$ is,
In the cluster centre, it is very likely that two (or even more)
stars come very close together in a hyperbolic encounter.
%and their relative distance becomes smaller than $p_{90}$.
%Energetic binaries may form via three-body encounters.
%, if there are not already primordial binaries.
%
As the relative distance of the two bodies becomes small ($R \rightarrow 0$),
their timesteps are reduced to prohibitively small values, and
truncation errors grow due to the singularity in the gravitational
potential, eqs.~(\ref{eq:accel0}) and (\ref{eq:accel1}).
In the NBODY code, the parameter RMIN is used to define a close
encounter, and it is kept to the value of equation \ref{eq:p90}
(if KZ(16) > 0 is chosen in the control parameters).
%%%/* But the user can specify it as an input parameter freely (if KZ(16)=0). */
%It is a good idea to keep it close to $p_{90}$, but depending on the
%astrophysical system simulated one may want to deviate even strongly 
%from this.
The corresponding time step DTMIN can be estimated from 
%
\begin{equation}
dt_{\rm min} = \kappa \Bigl[{\eta \over 0.03}\Bigr]
       \Bigl({r_{\rm min}^3 \over \langle m \rangle }\Bigr)^{1/2}
          \label{eq:dtmin}
\end{equation}
%
where $\kappa$ is a free numerical factor,
$\eta$ the general time step factor,
and $\langle m \rangle$ the average stellar mass \cite{a_85}.
%The code uses DTMIN as an input parameter, so the user can override
%the value of equation \ref{eq:dtmin} again via KZ(16)=0.
If two particles are getting closer to each other than RMIN, and
their time steps getting smaller than DTMIN, then they are candidates
for ``regularization''.

Regularization is an elegant trick in order to deal with such particles
which are as close as the diamond in the Figure \ref{pic:neighbour}.
The idea is to take both stars out of the main integration cycle,
replace them by their centre of mass (c.m.) and advance the
usual integration with this composite particle instead of
resolving the two components.
The two members of the regularized pair (henceforth KS pair)
will be relocated to the beginning of all vectors containing
particle data, while at the end one additional c.m.\ particle is
created (see below).
One of the purposes of the code variable NAME(I) is to identify
particles after such a reshuffling of data.

To be actually regularized, the two particles have to fulfil two
more sufficient criteria: that they are approaching each other,
and that their mutual force is dominant.
In the equations in routine {\tt search.f}, these sufficient criteria
are defined as
%
\begin{eqnarray}
& {\mathbf R}\cdot {\mathbf V} > 0.1 \sqrt{(G(m_1+m_2)R} \nonumber\\
& \gamma := {\vert {\mathbf a}_{\rm pert}\vert \cdot R^2 \over G(m_1+m_2) } < 0.25 \nonumber
%\label{eq:suff}
\end{eqnarray}
%
Here, ${\mathbf a}_{\rm pert}$ is the vectorial differential force
exerted by other perturbing particles onto the two candidates,
$R$, ${\mathbf R}$, ${\mathbf V}$ are scalar and vectorial distance
and relative velocity vector between the two candidate, respectively.
The factor 0.1 in the upper equation allows nearly circular
orbits to be regularized;
$\gamma < 0.25$ demands that the relative strength of the perturbing
forces to the pairwise force is one quarter of the maximum.
These conditions describe quantitatively that a two-body subsystem
is dynamically separated from the rest of the system, but not
unperturbed.

The internal motion of a KS pair will be determined by switching
to a different (regularized) coordinate system. % Aarseth (1993)
%The transformation employed is the Kustaanheimo--Stiefel
%transformation \cite{ks_65}.
This transformation can be traced back to the square in quaternion
space, where --- by sacrificing some commutativity rules --- it is
guaranteed that the real-space motion does not leave the
three-dimensional Cartesian space.
It involves a set of four regular spatial coordinates and a
fictitious time $s(t)$, obtained in its simplest variant by the
transformation $dt = R ds$. % Aarseth (1985)
Any unperturbed two--body orbit in real space is mapped onto
a harmonic oscillator in KS--space with double the frequency.
Since the harmonic potential is regular, numerical integration with
high accuracy can proceed with much better efficiency, and there
is no danger of truncation errors for arbitrarily small separations.
The internal time--step of such a KS--regularized pair is independent
of the eccentricity and, depending on the parameter ETAU, of the
order of some 50--100 steps per orbit. 
% Aarseth (1993), p14, advantages
%
The method of regularization goes back to \cite{ks_65} and makes an
accurate calculation of a perturbed two--body motion possible.
A modern theoretical approach to this subject can be found in
\cite{ns_92};
the Hamiltonian formalism of the underlying transformations is
nicely explained in \cite{m_97}.

While regularization can be used for any analytical two--body
solution across a mathematical collision, it is practically applied
to perturbed pairs only.
Once the perturbation $\gamma$ falls below a critical value
(input parameter GMIN $\approx 10^{-6}$),
a KS--pair is considered unperturbed, and the analytical solution for
the Keplerian orbit is used instead of doing numerical integration.
A little bit misleading is that such unperturbed KS--pairs are
denoted in the code as ''mergers'', e.g. in the number or merges (NM)
and the energy of the mergers (EMERGE).
Merged pairs can be resolved at any time if the perturbation changes.
The two--body KS regularization occurs in the code either for
short-lived hyperbolic encounters or for persistent binaries. 

In the code, the KS--pair appears as a new particle at the postion
of the centre of mass.
The variable NTOT, that contains the total number of particles $N$
plus the c.m.'s, is increased by 1.
When the pair is disrupted, NTOT is decreased again.
The maximum number of possible KS--pairs is saved in the variable
KMAX, which sets a threshold for the extension of the vector NTOT
(see Chapter~\ref{ch:thresholds}).\\

Close encounters between single particles and binary stars are
also a central feature of cluster dynamics.
Such temporary triple systems often reveal irregular motions, 
ranging from just a perturbed encounter to a very complex interaction,
in which disruption of binaries, exchange of components and ejection
of one star may occur.
Although not analytically solvable, the general three--body problem
has received much attention.
So, the KS--regularization was expanded to the isolated 3-- and 4--body
problem, and later on to the perturbed 3--, 4--, and finally to the
$N$--body problem.
The routines are called
%
\begin{itemize}
\item {\tt triple.f}$\:$ (unperturbed 3-body subsystems, \cite{az_74}),
\item {\tt quad.f}$\:$  (unperturbed 4-body subsystems), and
\item {\tt chain.f}$\:$ with different stages of implementation
      (slow-down, Stumpff functions, see for consecutive references
      Mikkola \& Aarseth 1990, 1993, 1996, 1998, and \cite{m_97}).
\end{itemize}
%
While occurrences of ``triple'' and ``quad'' will be rare in a
simulation, the chain regularization is invoked if a KS--pair
has a  close encounter with another single star or another pair.
Especially, if systems start with a large number of primordial
binaries, such encounters may lead to stable (or quasi-stable)
hierarchical triples, quadruples, and higher multiples.
They have to be treated by using special stability criteria.
Some of them are actually already implemented, but there is ongoing
research and development in the field.
%in particular (as relevant for our codes) by Rosemary Mardling and Sverre Aarseth.

A typical way to treat all such special higher subsystems is to define
their c.m.\ to be a pseudo-particle, i.e.\ a particle with a
known sub-structure (very much like nodes in a TREE code).
The members of the pseudo-particles will be deactivated by setting
their mass to zero (ghost particles).
At present there can only exist one chain at a time in the code,
while merged KS binaries, and hierarchical subsystems can be more
frequent.
Details of these procedures are beyond the scope of this introductory
manual.

Every subsystem --- KS pair, chain or hierarchical subsystem ---
is perturbed.
Perturbers are typically those objects that get closer to the object
than $R_{\rm sep} = R/\gamma_{\rm min}^{1/3}$, where $R$ is the
typical size of the subsystem;
for perturbers, the components of the subsystem are resolved in their
own force computation as well (routines {\tt cmfreg.f}, {\tt cmfirr.f}).


